\section{范围、架构与可审计边界}
\label{sec:ts-scope}

\subsection{问题定义}
给定富语义交直流系统 $\mathcal S$、离散 profile 集合
$\{\pi_k(t)\}_{t=0}^{T-1}$、步长 $\Delta t$ 和设备参数，模块计算
\[
  \{u_{g,t},p_{g,t},e_{s,t},V_t,I_t\}_{t=0}^{T-1},
\]
并将功率积分为 MWh、将可计价项积分为货币量。入口按保真度分成三层：
\begin{enumerate}
  \item \textbf{时序层}：\srcpath{solve\_time\_series\_pf}，单个耦合窗口内执行
        UC（可选）、逐步 AC-OPF（可选）和 PF 校验；
  \item \textbf{年度层}：\srcpath{solve\_annual\_production\_simulation}，提供
        L0 年度块、L2 周 UC、L3 日 replay，或绕过层级的独立按日并行；
  \item \textbf{生命周期层}：\srcpath{run\_lifecycle\_simulation}，年循环更新负荷、PV
        降额和储能 SOH，再调用年度 schedule-only 快速层；\srcpath{run\_lifecycle\_comparison}
        在克隆系统上扫描容量倍率。
\end{enumerate}

\subsection{数据流与身份}
所有求解均遵循“rich system $\rightarrow$ profile snapshot $\rightarrow$ canonical projection
$\rightarrow$ solver result $\rightarrow$ rich attribution”的主线。设备的稳定
\fld{index}、向量位置和图节点编号绝不混用；\fld{UCSchedule} 的行是\textbf{在役设备的自然顺序}
（不是稳定 ID 排序），而 \fld{opf\_results}/\fld{pf\_results} 的外层索引是局部窗口的 $t$。
年度拼接时由 \fld{global\_step=start\_step+t} 恢复全局时标。AC 与 DC 母线即使具有相同整数
ID 也使用不同域映射。

\subsection{两条年度路径}
\begin{center}\footnotesize
\begin{tabular}{@{}L{3.3cm}L{5.2cm}L{5.8cm}@{}}
\toprule
路径 & 实际执行 & 主要假设\\ \midrule
层级 sequential & L0 heuristic $\rightarrow$ 周 UC（168 h + lookahead）$\rightarrow$ 日 replay & 周间 SOC 重新锚定；L0 预算/维护目前是默认占位值\\
按日 parallel & DynamicOPF 日窗独立求解；固定式 AC/DC 储能末端 SOC 强制回到初值 & 日与日之间没有能量套利；UC/移动储能状态不准入\\
\bottomrule
\end{tabular}
\end{center}

\begin{gapnote}
\textbf{已实现：}日窗划分、DynamicOPF 年度准入、固定储能逐日循环 SOC、结果按全局步拼接、
并行后端守卫和 provenance 字段。SCUC/DynamicSCED 的日配置只保留为内部结构，年度入口不准入。
\textbf{仅结构存在：}L0 维护、燃料预算和
\fld{iterative\_feedback} 参数；当前 \srcpath{solve\_annual\_plan} 不求解 L0 MILP，
\fld{max\_feedback\_iterations} 与 \fld{budget\_violation\_tol\_mwh} 也未参与控制流。
\end{gapnote}

\subsection{复杂度与内存}
设 $G,S,R$ 分别为在役机组、储能和可再生数，日窗长度为 $\tau$。单个 UC 窗口的二进制量
主要为 $O(G\tau)$，年度层级将其变成约 52 个周问题或 $\lceil T/\tau\rceil$ 个日问题；
\fld{step\_results} 为 $O(T)$，PF 快照只在 \fld{pf\_snapshot\_interval} 命中时保留。
并行日路径按日结果槽位写入，随后主线程按日序拼接，因此数值顺序不依赖线程完成顺序。
